% =====================================================================
\chapter{Codificación aritmética: predecir es comprimir}
\label{tema:4}
% =====================================================================

Este tema presenta la codificación aritmética, que representa un mensaje completo mediante un subintervalo de $[0,1)$ de anchura igual a su probabilidad, y demuestra que con ella cualquier modelo probabilístico se convierte en un compresor cuya longitud es, salvo dos bits, la log-verosimilitud negativa del mensaje en bits. La construcción sigue la codificación de Shannon--Fano--Elias de~\cite[§5.9]{ct} y su extensión a secuencias de~\cite[§13.3]{ct}; la presentación como \emph{stream code} con un modelo adaptativo es la de MacKay~\cite[cap.~6]{mackay}, y la implementación práctica, la de Witten, Neal y Cleary~\cite{wnc}. La conexión con los modelos de lenguaje sigue a Delétang \emph{et al.}~\cite{deletang} y a Gao \emph{et al.}~\cite{pile}.

\section{Intervalos y probabilidades}

Sea $\X=\{x^{(1)},\dots,x^{(K)}\}$ un alfabeto con un orden fijado y sea $q$ un modelo autorregresivo, que asigna a cada contexto $x_{1:t-1}$ una distribución $q(\cdot\mid x_{1:t-1})$ sobre $\X$. La \emph{probabilidad acumulada} de $x$ en la posición $t$ es la masa de los símbolos anteriores en el orden fijado:
\[
F_t(x)=\sum_{x'<x}q(x'\mid x_{1:t-1}).
\]

\begin{definicion}[Intervalo de una secuencia]\label{def:intervalo}
El \emph{intervalo} de la secuencia $x_{1:n}$ es $I(x_{1:n})=[a_n,a_n+w_n)$, donde $a_0=0$, $w_0=1$ y, para $t=1,\dots,n$,
\[
a_t=a_{t-1}+w_{t-1}\,F_t(x_t),\qquad w_t=w_{t-1}\,q(x_t\mid x_{1:t-1}).
\]
\end{definicion}

Cada paso divide el intervalo actual en $K$ subintervalos consecutivos, de anchuras proporcionales a las probabilidades del símbolo siguiente, y se queda con el del símbolo observado.

\begin{proposicion}\label{prop:intervalos}
Para todo $n\ge 0$:
\begin{enumerate}
  \item la anchura de $I(x_{1:n})$ es $w_n=q(x_{1:n})=\prod_{t=1}^{n}q(x_t\mid x_{1:t-1})$;
  \item los intervalos de todas las secuencias de longitud $n$ forman una partición de $[0,1)$;
  \item $I(x_{1:n})\subseteq I(x_{1:n-1})$ si $n\ge 1$.
\end{enumerate}
\end{proposicion}

\begin{proof}
(i) Por inducción, $w_n=w_{n-1}\,q(x_n\mid x_{1:n-1})$ y la regla de la cadena del módulo~1 da el producto. (ii) y (iii) Por inducción sobre $n$: para $n=0$ hay un único intervalo, $[0,1)$. Si los intervalos de longitud $n-1$ forman una partición, cada uno, $[a,a+w)$, se divide en los subintervalos $[a+wF_n(x),\,a+wF_n(x)+w\,q(x\mid x_{1:n-1}))$, que son consecutivos porque $F_n(x)+q(x\mid x_{1:n-1})$ es la acumulada del símbolo siguiente a $x$, y que cubren $[a,a+w)$ porque las probabilidades suman~1.
\end{proof}

\begin{ejemplo}\label{ej:aab}
Sea $\X=\{\mathrm{a},\mathrm{b},\mathrm{c}\}$, en ese orden, con el modelo sin memoria $q=(0{,}6;\ 0{,}3;\ 0{,}1)$ y acumuladas $F=(0;\ 0{,}6;\ 0{,}9)$. Para el mensaje «aab»:
\begin{center}
\begin{tabular}{@{}cccc@{}}
\toprule
Paso & Símbolo & Intervalo & Anchura \\
\midrule
0 & --- & $[0;\ 1)$ & $1$ \\
1 & a & $[0;\ 0{,}6)$ & $0{,}6$ \\
2 & a & $[0;\ 0{,}36)$ & $0{,}36$ \\
3 & b & $[0{,}216;\ 0{,}324)$ & $0{,}108$ \\
\bottomrule
\end{tabular}
\end{center}
La anchura final es $0{,}6\cdot 0{,}6\cdot 0{,}3=0{,}108=q(\text{aab})$, y $-\log 0{,}108=3{,}211$ bits.
\end{ejemplo}

\begin{figure}[h]
\centering
\begin{tikzpicture}[x=11cm,y=1cm,font=\small]
  % nivel 1
  \draw[thick] (0,2) -- (1,2);
  \foreach \p/\l in {0/0,0.6/{0{,}6},0.9/{0{,}9},1/1} \draw (\p,1.9) -- (\p,2.1) node[above] {$\l$};
  \node at (0.3,1.7) {a}; \node at (0.75,1.7) {b}; \node at (0.95,1.7) {c};
  % nivel 2
  \draw[thick] (0,1) -- (0.6,1);
  \foreach \p in {0,0.36,0.54,0.6} \draw (\p,0.9) -- (\p,1.1);
  \node at (0.18,0.7) {aa}; \node at (0.45,0.7) {ab}; \node at (0.57,0.7) {ac};
  % nivel 3
  \draw[thick] (0,0) -- (0.36,0);
  \foreach \p in {0,0.216,0.324,0.36} \draw (\p,-0.1) -- (\p,0.1);
  \node at (0.108,-0.3) {aaa}; \node at (0.27,-0.3) {aab}; \node at (0.342,-0.3) {aac};
  \draw[line width=3pt,regla] (0.216,0) -- (0.324,0);
  \draw[dashed,regla] (0,1.9) -- (0,0.1) (0.6,1.9) -- (0.6,1.1) (0.36,0.9) -- (0.36,0.1);
  \node[below] at (0.216,-0.45) {$0{,}216$}; \node[below] at (0.324,-0.45) {$0{,}324$};
\end{tikzpicture}
\caption{Codificación aritmética del mensaje «aab» (\cref{ej:aab}). Cada nivel subdivide el intervalo del nivel anterior en proporción a las probabilidades del símbolo siguiente.}
\label{fig:aritmetica}
\end{figure}

\section{Palabras de código y longitud}

Para transmitir el intervalo basta enviar una cadena binaria $c$ cuyo intervalo diádico $J(c)$ (\cref{tema:3}) esté contenido en $I(x_{1:n})$: el receptor sabrá entonces que el número enviado, cualesquiera que sean los bits que lo sigan, está dentro del intervalo del mensaje.

\begin{teorema}\label{teo:palabra}
Sea $I=[a,a+w)\subseteq[0,1)$ con $w>0$, y sea $\ell=\lceil-\log w\rceil+1$. Existe una cadena binaria $c$ de longitud $\ell$ con $J(c)\subseteq I$.
\end{teorema}

\begin{proof}
Por la elección de $\ell$, $2^{-\ell}\le w/2$. Sea $m=\lceil a\,2^{\ell}\rceil$, de modo que $m2^{-\ell}$ es el menor múltiplo de $2^{-\ell}$ que no es menor que $a$; entonces $a\le m2^{-\ell}<a+2^{-\ell}$, y
\[
(m+1)\,2^{-\ell}<a+2\cdot2^{-\ell}\le a+w.
\]
Así, $[m2^{-\ell},(m+1)2^{-\ell})\subseteq I$. Como $(m+1)2^{-\ell}\le 1$, se tiene $0\le m<2^{\ell}$, y $c$ es la escritura binaria de $m$ con $\ell$ cifras.
\end{proof}

\begin{corolario}[Longitud de la codificación aritmética]\label{cor:longitud}
Para cada $n$, asignar a cada secuencia $x_{1:n}$ con $q(x_{1:n})>0$ la cadena del \cref{teo:palabra} para $I(x_{1:n})$ define un código prefijo sobre las secuencias de longitud $n$, con
\[
\ell(x_{1:n})=\lceil-\log q(x_{1:n})\rceil+1<-\log q(x_{1:n})+2.
\]
\end{corolario}

\begin{proof}
Los intervalos de secuencias distintas de longitud $n$ son disjuntos (\cref{prop:intervalos}), luego también lo son los intervalos diádicos contenidos en ellos, y ninguna palabra es prefijo de otra. La cota de la longitud se sigue de $\lceil y\rceil<y+1$ y de $w_n=q(x_{1:n})$.
\end{proof}

\begin{ejemplo}\label{ej:aab-palabra}
Para «aab», $\ell=\lceil 3{,}211\rceil+1=5$ y $m=\lceil 0{,}216\cdot 32\rceil=\lceil 6{,}912\rceil=7$, que da la palabra \bin{00111}, con $J(\bin{00111})=[0{,}21875;\ 0{,}25)\subseteq[0{,}216;\ 0{,}324)$. La cota es una garantía, no siempre el óptimo: la cadena \bin{0100}, con $J(\bin{0100})=[0{,}25;\ 0{,}3125)$, también cabe en el intervalo y solo tiene 4 bits; ninguna de 3 bits cabe, porque los intervalos diádicos de anchura $0{,}125$ son $[0{,}125;\ 0{,}25)$ y $[0{,}25;\ 0{,}375)$, y ninguno está contenido en $[0{,}216;\ 0{,}324)$.
\end{ejemplo}

\paragraph{Descodificación.} El receptor conoce el modelo $q$ y la longitud $n$ (o bien el alfabeto incluye un símbolo de fin, como el símbolo \fin\ del módulo~1). Recibida $c$, para $t=1,\dots,n$ subdivide el intervalo actual según $q(\cdot\mid x_{1:t-1})$, determina el único símbolo $x_t$ cuyo subintervalo contiene a $J(c)$ y lo añade al contexto. Por la \cref{prop:intervalos}, ese símbolo existe, es único y coincide con el enviado. En la práctica, la aritmética de precisión infinita se sustituye por aritmética entera con renormalización, que emite los bits a medida que se fijan, con un sobrecoste despreciable~\cite{wnc}.

\section{Predecir es comprimir}

\begin{teorema}[Compresión con un modelo]\label{teo:compresion}
Sea $q$ un modelo autorregresivo cualquiera. La codificación aritmética con $q$ codifica cada secuencia $x_{1:n}$ con
\[
\ell(x_{1:n})<-\log q(x_{1:n})+2=\sum_{t=1}^{n}-\log q(x_t\mid x_{1:t-1})+2\quad\text{bits}.
\]
Si $X_{1:n}$ sigue una distribución $p$, la longitud media cumple
\[
H(X_{1:n})\le\E[\ell(X_{1:n})]<H(p,q)+2=H(X_{1:n})+\KL{p}{q}+2,
\]
donde $p$ y $q$ denotan aquí las distribuciones de la secuencia completa.
\end{teorema}

\begin{proof}
La primera desigualdad es el \cref{cor:longitud} junto con la regla de la cadena. Tomando esperanzas con $X_{1:n}\sim p$, $\E[-\log q(X_{1:n})]=H(p,q)$, que es igual a $H(X_{1:n})+\KL{p}{q}$ por la \cref{prop:descomposicion}. La cota inferior es el \cref{teo:cota-inferior}, aplicado al código prefijo sobre secuencias de longitud $n$.
\end{proof}

\begin{corolario}\label{cor:por-simbolo}
Si las $X_t$ son independientes con distribución $p$ y el modelo es sin memoria con distribución $q$, la longitud media por símbolo cumple
\[
H(p)\le\frac{\E[\ell(X_{1:n})]}{n}<H(p)+\KL{p}{q}+\frac{2}{n}.
\]
Con $q=p$, la codificación aritmética alcanza la entropía con un sobrecoste de $2/n$ bits por símbolo, sin tablas que crezcan con $n$.
\end{corolario}

\begin{proof}
Por independencia, $H(X_{1:n})=nH(p)$ y $\KL{p^{\otimes n}}{q^{\otimes n}}=n\KL{p}{q}$, porque el logaritmo del cociente de productos es la suma de los logaritmos. Basta dividir el \cref{teo:compresion} por $n$.
\end{proof}

El \cref{teo:compresion} es la mitad de la equivalencia anunciada: \emph{todo predictor da un compresor}, y el tamaño del archivo comprimido es la suma de las sorpresas del modelo. La otra mitad se vio en el \cref{teo:cota-inferior}: \emph{todo compresor da un predictor}, porque las longitudes de un código prefijo definen la distribución $r(x)\propto2^{-\ell(x)}$. Mejorar uno equivale a mejorar el otro.

\begin{figure}[h]
\centering
\begin{tikzpicture}[font=\small,node distance=5mm and 9mm,
  caja/.style={draw,thick,rounded corners=2pt,minimum height=8mm,align=center,inner sep=4pt},
  >={Stealth[length=2mm]}]
  \node[caja] (txt) {texto\\$x_{1:n}$};
  \node[caja,right=of txt] (cod) {codificador\\aritmético};
  \node[caja,above=of cod] (m1) {modelo $q$};
  \node[caja,right=of cod] (bits) {bits\\$\approx-\log q(x_{1:n})$};
  \node[caja,right=of bits] (dec) {descodificador\\aritmético};
  \node[caja,above=of dec] (m2) {mismo\\modelo $q$};
  \node[caja,right=of dec] (txt2) {texto\\$x_{1:n}$};
  \draw[->] (txt) -- (cod); \draw[->] (m1) -- (cod); \draw[->] (cod) -- (bits);
  \draw[->] (bits) -- (dec); \draw[->] (m2) -- (dec); \draw[->] (dec) -- (txt2);
  \draw[->,dashed] (dec.south) |- ++(0,-0.6) -| node[pos=0.25,below] {token recuperado como contexto} ([xshift=6mm]m2.east) -- (m2.east);
\end{tikzpicture}
\caption{Compresión con un modelo de lenguaje. Emisor y receptor usan el mismo modelo; el receptor recupera cada token y lo añade al contexto para obtener la distribución del siguiente.}
\label{fig:compresor}
\end{figure}

\begin{observacion}[Conexión con los LLM]
La pérdida de entrenamiento de un LLM, expresada en bits y multiplicada por el número de tokens, es la longitud del texto comprimido con ese modelo por codificación aritmética, salvo dos bits. Minimizar la pérdida es, literalmente, minimizar el tamaño de un archivo. Delétang \emph{et al.}~\cite{deletang} usaron modelos de lenguaje como compresores y mostraron que Chinchilla~70B, entrenado sobre todo con texto, comprime fragmentos de imágenes de ImageNet al $43{,}4\,\%$ de su tamaño y muestras de audio de LibriSpeech al $16{,}4\,\%$, mejor que los compresores específicos PNG ($58{,}5\,\%$) y FLAC ($30{,}3\,\%$). Los mismos autores señalan que, si se cuenta en el tamaño comprimido también el del propio modelo, el tamaño del conjunto de datos limita el tamaño de modelo que compensa usar.
\end{observacion}

\section{Bits por byte}

La pérdida por token y la perplejidad dependen del tokenizador: un vocabulario mayor produce menos tokens, cada uno más difícil de predecir, y la cifra por token cambia aunque el modelo comprima igual de bien. Para comparar modelos con tokenizadores distintos se normaliza por la longitud del texto, que no depende del tokenizador.

\begin{definicion}[Bits por byte]\label{def:bpb}
Sea un texto de $B$ bytes en la codificación UTF-8, dividido por el tokenizador en $N$ tokens. Los \emph{bits por byte} del modelo $q$ sobre el texto son
\[
\bpb=\frac{1}{B}\sum_{t=1}^{N}-\log_2 q(x_t\mid x_{1:t-1}).
\]
\end{definicion}

\begin{proposicion}\label{prop:bpb}
Si $\mathcal{J}$ es la pérdida media en nats por token y $\bar b=B/N$ el número medio de bytes por token, entonces $\bpb=\mathcal{J}/(\bar b\ln 2)$.
\end{proposicion}

\begin{proof}
$\sum_t-\log_2q(x_t\mid x_{1:t-1})=\frac{1}{\ln 2}\sum_t-\ln q(x_t\mid x_{1:t-1})=\frac{N\mathcal{J}}{\ln 2}$. Dividiendo por $B=N\bar b$ se obtiene el resultado.
\end{proof}

Por el \cref{teo:compresion}, el bpb es, salvo un término despreciable, el tamaño del texto comprimido con el modelo dividido por su tamaño original en bytes; un bpb de $b$ significa que el texto se comprime a una fracción $b/8$ de su tamaño.

\begin{ejemplo}\label{ej:bpb}
Un modelo con pérdida $\mathcal{J}=2{,}0$ nats por token y un tokenizador de $4$ bytes por token tiene $\bpb=2{,}0/(4\cdot 0{,}6931)=0{,}721$ bits por byte: comprime el texto a $0{,}721/8\approx 9{,}0\,\%$ de su tamaño.
\end{ejemplo}

\begin{center}
\begin{tabular}{@{}llcl@{}}
\toprule
Métrica & Unidad & Tokenizador & Uso habitual \\
\midrule
pérdida & nats por token & depende & entrenamiento \\
perplejidad & $e^{\text{pérdida}}$ & depende & número efectivo de opciones \\
bits por carácter & bits por carácter & no depende & modelos de caracteres \\
bits por byte & bits por byte & no depende & comparar modelos distintos \\
\bottomrule
\end{tabular}
\end{center}

Los bits por carácter tampoco dependen del tokenizador, pero sí de cómo se cuentan los caracteres; los bytes UTF-8 están definidos sin ambigüedad. En español, las vocales con tilde, la \emph{ñ} y las comillas angulares ocupan 2 bytes cada una, de modo que un mismo texto tiene más bytes que caracteres y su bpb es algo menor que sus bits por carácter. El artículo del conjunto de datos The Pile~\cite{pile} popularizó el bpb para comparar modelos con tokenizadores distintos. Anthropic, Google y xAI no publican el bpb de sus modelos actuales.

\begin{lecturas}
\item Cover y Thomas~\cite{ct}, §5.9 «Shannon--Fano--Elias Coding» y §13.3 «Arithmetic Coding».
\item MacKay~\cite{mackay}, cap.~6 «Stream Codes».
\item Witten, Neal y Cleary~\cite{wnc}, implementación clásica con aritmética entera.
\item Delétang \emph{et al.}~\cite{deletang}, \emph{Language Modeling Is Compression}.
\item Gao \emph{et al.}~\cite{pile}, uso de los bits por byte como métrica de evaluación.
\end{lecturas}

\section*{Ejercicios del tema 4}
\addcontentsline{toc}{section}{Ejercicios del tema 4}

\begin{ejercicio}{1}{eje:4-bpb-concepto}
Explique por qué se comparan modelos de lenguaje con tokenizadores distintos en bits por byte y no en perplejidad por token. Ilústrelo con dos modelos: A gasta 2 bits por token con 2 bytes por token y B gasta 4 bits por token con 4 bytes por token.
\end{ejercicio}

\begin{ejercicio}{2}{eje:4-bac}
Con el modelo sin memoria $q(\mathrm{a})=1/2$, $q(\mathrm{b})=1/4$, $q(\mathrm{c})=1/4$, en ese orden, codifique aritméticamente el mensaje «bac». Dé el intervalo final, la longitud ideal $-\log q(\text{bac})$, la palabra de código más corta cuyo intervalo diádico esté contenido en el intervalo final y la longitud que garantiza el \cref{cor:longitud}.
\end{ejercicio}

\begin{ejercicio}{2}{eje:4-descodificar}
Con el mismo modelo del \cref{eje:4-bac}, un receptor sabe que el mensaje tiene 3 símbolos y recibe la palabra \bin{01110}. Descodifique el mensaje paso a paso y compruebe que la longitud recibida coincide con $-\log q$ del mensaje.
\end{ejercicio}

\begin{ejercicio}{2}{eje:4-abbb}
Una fuente binaria emite A con probabilidad $1/4$ y B con probabilidad $3/4$. Codifique el mensaje «ABBB» con el código de Huffman símbolo a símbolo y con codificación aritmética (orden A, B), dando en este caso la palabra más corta cuyo intervalo diádico cabe en el intervalo final. Compare ambas longitudes con la longitud ideal y comente qué ocurre con mensajes largos.
\fuente{adaptado de Universidad de Cambridge, \emph{Information Theory}~\cite{camit}, \emph{Examples Sheet}, \emph{Supervision}~2, ejercicio~5 (el original usa un mensaje de nueve símbolos e incluye Lempel--Ziv).}
\end{ejercicio}

\begin{ejercicio}{2}{eje:4-bpb}
Un modelo alcanza una pérdida de $0{,}7$ nats por token con un tokenizador de $3{,}5$ bytes por token. Calcule sus bits por byte y su factor de compresión respecto de 8 bits por byte.
\end{ejercicio}

\begin{ejercicio}{2}{eje:4-equivocado}
Los datos siguen $p=(1/4,1/4,1/2)$ sobre $\{\mathrm{a},\mathrm{b},\mathrm{c}\}$, pero el código se ha diseñado para $q=(1/2,1/4,1/4)$ con longitudes $\ell(x)=-\log q(x)$. Calcule la longitud media del código bajo $p$, $H(p)$ y $\KL{p}{q}$, y compruebe que la longitud media es $H(p)+\KL{p}{q}$.
\end{ejercicio}

\begin{ejercicio}{3}{eje:4-equivocado-teo}
Sea un código de Shannon diseñado para $q$, con $\ell(x)=\lceil-\log q(x)\rceil$, usado con datos que siguen $p$ (con $q(x)>0$ siempre que $p(x)>0$). Demuestre que
\[
H(p)+\KL{p}{q}\le\E_{p}[\ell(X)]<H(p)+\KL{p}{q}+1.
\]
\fuente{adaptado de Cover y Thomas~\cite{ct}, §5.4, teorema sobre la codificación con una distribución equivocada.}
\end{ejercicio}
